\section{Circuits, counting, alternation, and finite probability}
\label{sec:extensions}

These extensions illustrate that the foundation is reusable. They involve
finite sums, products and syntax, with resource bounds proved separately.
They require neither infinite probability spaces nor quantification over an
actual infinite family as a new sort.

\subsection{Uniform and nonuniform circuits}

For a Boolean presentation \(A(x)\), a polynomial-size circuit presentation
means that for each length \(n\) there exists a circuit \(C_n\) of at most
\(p(n)\) gates agreeing with \(A\) on all \(2^n\) words of that length.
This is the ordinary nonuniform class P/poly, read schematically over \(A\)
or arithmetically when \(A\) is a supplied Boolean numerical function.
There is no automatic efficient operation that selects \(C_n\) from an
existence statement.

If a fixed polynomial-time machine computes \(A\), unroll its computation
with all \(n\) input bits left as circuit inputs, instead of hardwiring
\(x\) as in Cook--Levin. The same cell induction proves correctness, and
the number of gates is polynomial in \(n\). Moreover a generator constructs
\(C_n\) in polynomial time from \(1^n\). Conversely, if a polynomial-time
generator on \(1^n\) produces correct polynomial-size circuits, generating
and evaluating the circuit is a polynomial-time decision procedure. Thus
\ipra{} supports both the nonuniform upper-bound construction and the
equivalence with this particular polynomial-time uniformity convention.
The unary presentation of \(n\) is part of that convention.

\begin{theorem}[Constructively selected hard truth tables]
\label{thm:circuitcount}
There is a primitive-recursive function \(H\) such that, for every \(n\ge4\),
\(H(n)<2^{2^n}\) is the truth table of a Boolean function not computed by
any \(n\)-input circuit of at most
\[
 s(n)=\left\lfloor\frac{2^n}{16n}\right\rfloor
\]
internal AND, OR, or NOT gates. Constants and inputs are available as wires.
For each such circuit a primitive-recursive bounded search returns an input
on which it disagrees with this truth table. These are \ipra{} theorems.
\end{theorem}

\begin{proof}
Put \(s=s(n)\) and \(R=n+s+2\), counting two constant wires. Pad a circuit
with unused gates to exactly \(s\) gates; this preserves its function.
An overcount of the exactly-\(s\) descriptions is
\[
 K=R(3R^2)^s:
\]
choose an output wire, and for each gate one of three operations and two
predecessor indices. Ignore invalid forward references; doing so only
increases the count. NOT gates simply ignore their second index. Thus
\(K\) bounds the number of \emph{functions} represented using at most
\(s\) gates. We do not claim that padding is injective on all shorter
syntactic descriptions.

For \(n\ge4\), \(n+2\le2^{n-1}\) and \(s\le2^{n-1}\), hence
\(R\le2^n\). Also \(3R^2\le2^{4n}\), so
\[
 K\le2^{n+4ns},\qquad
 n+4ns\le n+2^{n-2}\le2^{n-1}<2^n.
\]
The inequality \(n\le2^{n-2}\) starts at \(n=4\) and is preserved by
successor; the other inequalities follow by multiplication and integer
division. Hence \(K<2^{2^n}\), the number of possible truth tables.

Enumerate the \(K\) candidate descriptions, mapping invalid ones to a
fixed table, and compute each valid circuit's \(2^n\) output bits. Maintain
an image bitset indexed by the \(2^{2^n}\) possible truth tables. After
processing \(i\) descriptions, its number of set bits is at most \(i\):
inserting one value adds at most one new mark. This invariant uses bounded
sums and Boolean comparisons, so quantifier-free induction proves it.
Since the image cannot cover the table universe, bounded least search
returns a missing table. Define this table to be \(H(n)\), taking a fixed
default value for \(n<4\). Every loop and the image code itself have an
explicit primitive-recursive bound, even though they are enormous.

For a candidate circuit, equality with \(H(n)\) at all inputs would make
its truth-table integer equal to \(H(n)\), contrary to missingness.
More constructively, scan the \(2^n\) positions for the first differing
bit. If the scan had no success, its finite all-equal bit would be 1 and
the two table integers would be equal. Thus the scan supplies a mismatch
without turning an unbounded negation into an existence claim.
\end{proof}

This produces a primitive-recursive hard family, and therefore a
primitive-recursive language obtained by consulting the appropriate table
bit. Its construction can take far more than exponential time. It is not
an exponential circuit lower bound for NP, P, or another small uniform
complexity class. The example separates constructive existence from
efficient uniform construction with no ambiguity.

\subsection{Counting witnesses}

For a fixed polynomially bounded verifier, define
\[
 \#V(x)=\sum_{j<2^{p(n)}}
       \Acc_{\#}(e,x,\operatorname{bits}_{p(n)}(j),q(n)).
\]
This is a primitive-recursive natural-valued function and
\(0\le\#V(x)\le2^{p(n)}\). Its output has at most \(p(n)+1\) bits.
These presentations give the function class \(\#\mathsf P\).
Fixing exact certificate lengths makes the counted objects unambiguous;
padding conventions must be specified because they can change the count.

The support identity
\(\#V(x)>0\leftrightarrow N_v(x)=1\) follows from bounded witness search.
Disjoint tagged certificates add counts; pairs of certificates multiply
counts, by finite distributivity. Thus addition and multiplication of two
such functions have explicit verifier constructions after padding without
duplicating accepted certificates. For addition, use a tag and require all
padding bits to be zero; for multiplication, concatenate the two certificate
fields. This extra padding check is essential for exact counting.

When the circuit extension in Cook--Levin is unique for each input
assignment, counting satisfying gate assignments with the asserted output
counts exactly the accepting witnesses. This is a parsimonious statement
for that specified fixed-width witness representation. It does not assume
that an arbitrary many-one reduction preserves numbers of solutions.

\subsection{A fixed polynomial hierarchy}

For each fixed \(i\), consider presentations
\[
 \exists y_1<2^{p_1(n)}\ \forall y_2<2^{p_2(n)}\ \cdots\ Q_i y_i<2^{p_i(n)}
 \quad V(x,y_1,\ldots,y_i)=1,
\]
where the number of blocks \(i\) is fixed, each \(p_j\) is an explicit
polynomial, and \(V\) is a polynomially clocked Boolean verifier. These
are the bounded-quantifier presentations of \(\Sigma_i^P\); starting with
a universal block gives \(\Pi_i^P\). Every individual input formula is
bounded and therefore has a primitive-recursive characteristic bit.
Complement swaps the finite quantifiers and negates the matrix; this is
valid constructively because all the underlying finite tests are decidable.

A depth-first evaluation keeps one polynomial-length assignment for each
of the fixed number of blocks, the verifier workspace, and loop counters.
It therefore uses polynomial space. Its total number of branches is at
most \(2^{p_1(n)+\cdots+p_i(n)}\), with a polynomial factor for verification.
The resulting compiler proves that every fixed level lies in PSPACE.
Varying \(i\) as a metatheoretic family describes PH. This does not assert
that all levels collapse, nor that a single polynomial bound works for
arbitrary variable-level formulas. Section~\ref{sec:tqbf} handles a
variable-length quantified formula as its own finite input.

\subsection{Randomness is finite integer counting}

A randomized descriptor supplies a deterministic clocked machine, a
polynomial number \(p(n)\) of random bits, and a polynomial runtime.
Random bits are simply an additional exactly sized input field; early
termination ignores unused bits. For each input let
\[
 M=2^{p(n)},\qquad
 C(x)=\sum_{j<M}\operatorname{out}(x,\operatorname{bits}_{p(n)}(j)).
\]
The acceptance probability notation \(C(x)/M\) abbreviates a finite
rational number. Inequalities are represented by cross multiplication of
integers, so no measure theory or real numbers are needed.

There is a constructive subtlety when the intended language is an arbitrary
formula \(A(x)\). The two implications
\[
 A(x)\to 3C(x)\ge2M,\qquad
 \neg A(x)\to 3C(x)\le M
\]
alone need not identify \(A\) constructively: a high count gives only
\(\neg\neg A\). We therefore specify a BPP presentation by a proved gap
\[
 \forall x\ge1\quad
       (3C(x)\le M)\ \lor\ (3C(x)\ge2M)
\]
and pointwise equivalence of \(A(x)\) with the finite majority bit
\([2C(x)\ge M]\). Equivalently one can supply a Boolean characteristic
\(\chi_A\) and prove that the number of wrong-output seeds is at most
\(M/3\). The characteristic is not assumed to be computable in polynomial
time. In the standard extensional reading this is the usual BPP notion,
with its constructive identification made explicit.

\begin{theorem}[Integer error amplification]
\label{thm:amplification}
Suppose a randomized Boolean algorithm has at most \(F\) failing seeds out
of \(M\), where \(3F\le M\). The majority of
\(r=16(k+1)+1\) independent runs has at most a \(2^{-k}\) fraction of
failing combined seeds. The construction and the integer inequality are
provable in \ipra. For a polynomial-time randomized descriptor its runtime
and random-bit usage are polynomial in the input length and \(k\).
\end{theorem}

\begin{proof}
Write \(r=2m+1\) and consider all \(M^r\) tuples of seeds. Let \(S(\rho)\)
be the number of failing runs on tuple \(\rho\), and let \(F_0\le F\)
be the actual number of failing single-run seeds. Each successful seed
contributes weight 1 and each failing seed weight 2. Finite product
expansion gives
\[
 \sum_{\rho}2^{S(\rho)}=(M+F_0)^r\le(M+F)^r.
\]
If \(B\) tuples have a failing majority, each has weight at least
\(2^{m+1}\). Since \(3(M+F)\le4M\),
\[
 2^{m+1}B\le(M+F)^{2m+1},\qquad
 3\cdot9^m B\le2\cdot8^m M^{2m+1}.
\]
The second inequality follows by multiplying the first by
\(3^{2m+1}\), using the preceding bound, and canceling a positive common
power of two. All these are integer operations.

The concrete inequality \(2\cdot8^8=33554432<43046721=9^8\) gives, for
\(m=8(k+1)\),
\(2^{k+1}8^m<9^m\). Combining the inequalities yields
\(2^kB\le M^r\), the asserted bound. The proof of the product expansion
is induction over the number of independent seed fields; its invariant is
an equality of finite sums. Inequality preservation and counting bad
tuples are equally bounded. Thus only quantifier-free induction is used
after compilation.

Implement the independent runs sequentially, retaining a count of output
bits. The combined random input has \(r p(n)\) bits, and the runtime is
a fixed polynomial majorant of \(r\) times the one-run time plus input
and counter bookkeeping. This is polynomial in \(n,k\). No extraction of
a single good seed for all inputs, and no derandomization theorem, is
claimed.
\end{proof}

Deterministic machines ignore their random bits, giving
\(\Pclass\subseteq\mathsf{BPP}\). Exhaustive seed counting uses polynomial
space: retain one \(p(n)\)-bit seed, a \(p(n)+1\)-bit acceptance counter,
and the polynomial verifier workspace. Therefore
\(\mathsf{BPP}\subseteq\PSPACE\). A one-sided RP presentation has no
accepting seeds on no-instances and at least half its seeds accepting on
yes-instances, together with the same explicit language identification.
A successful seed is an NP witness, giving \(\mathsf{RP}\subseteq\NP\).
Complement gives \(\mathsf{coRP}\subseteq\coNP\). These arguments do not
assume that BPP equals P or that randomized expected-time bounds can always
be replaced by deterministic clocks with identical behavior.

\subsection{Worst case, finite averages, and amortized costs}

For a supplied total clocked algorithm, its actual stopping time below the
clock is a bounded least search. Maximum and minimum over the \(2^n\)
inputs are primitive recursive. For an explicitly given finite rational
distribution, multiply costs by its integer weights and divide by their
positive total weight. The average lies between the minimum and maximum
by summing the corresponding pointwise inequalities.

Amortized analysis concerns a different finite quantity. If states
\(s_0,\ldots,s_m\), nonnegative costs \(c_i\), and a numerical potential
\(\Phi\) obey
\(c_i+\Phi(s_{i+1})\le a_i+\Phi(s_i)\), summing and canceling gives
\[
 \sum_{i<m}c_i+\Phi(s_m)
       \le\sum_{i<m}a_i+\Phi(s_0).
\]
The induction is an inequality between finite sums. This provides the
usual amortized upper bound when \(\Phi\ge0\), and does not posit an
ordering of amortized and average-case complexity without additional
assumptions about workloads and distributions.
